% !TEX root = ../main.tex

\chapter{预期目标}

\section{预期目标}

本课题围绕Two-tone移位异或纠删码的高效实现，预期达到以下目标：

\begin{itemize}
    \item \textbf{算法实现目标}：完成Two-tone移位异或纠删码的高效编码与解码算法设计，实现支持多种丢失模式的统一解码流程。
    \item \textbf{系统集成目标}：将优化后的Two-tone算法以插件形式集成至Ceph分布式存储系统，实现功能性和性能测试。
    \item \textbf{性能优化目标}：通过缓存友好窗口划分、邻域合并访问等系统级优化，显著提升编码和解码吞吐量，降低内存访问开销。
    \item \textbf{对比验证目标}：与Ceph现有SOTA纠删码插件（如ISA-L、Cerasure、ZigZag）进行性能对比，验证本方案在实际应用中的优势。
    \item \textbf{理论与方法论目标}：总结移位异或纠删码的系统优化方法论，为后续相关算法研究和工程实现提供参考。
\end{itemize}